
Notons aussi que, @@M0@@ étant régulier donc normal, ses composantes connexes sont identiques à ses composantes irréductibles, lesquelles dominent @@M1@@, de sorte que l'homomorphisme @@M2@@ induit par l'inclusion @@M3@@ est bijectif. Par suite l'homomorphisme @@M4@@ induit par l'inclusion @@M5@@ peut aussi s'interpréter comme un homomorphisme
@@M6@@
qui mérite encore le nom de «~homomorphisme de spécialisation~». Ceci posé, on vérifie immédiatement sur les définitions que l'homomorphisme (7.7.1) est compatible avec les structures d'\emph{algèbres augmentées} sur @@M7@@ resp. @@M8@@ des deux membres définies dans (IV 2.8), compte tenu de l'homomorphisme précédent (7.7.3).

Supposant toujours @@M9@@ séparé, ce qui précède implique en particulier que (7.7.1) induit un homomorphisme d'anneaux gradués
@@M10@@
et utilisant la notion de classe de Chern de (V 6), on trouve un diagramme commutatif
@@M11@@

\subsection*{7.8. Spécialisation pour @@M12@@.}
Gardons les hypothèses de 7.7 : @@M13@@ séparé, @@M14@@ et @@M15@@ réguliers. Tenant compte de l'isomorphisme canonique (5.3.2) :
@@M16@@
l'homomorphisme (7.7.4) nous donne donc un homomorphisme canonique
@@M17@@
